\section{Asignación de actividades}

Sub equipos:
\begin{enumerate}
  \item Cortez Vega Jaime Fabian y Lujan De la Rosa Cristian Emir
  \item Arias Mendoza Angel Adrian y Vasquez Moreno Zair De Jesus
  \item Carrasco Pacheco Levi Josue y Nicolas Vasquez Joel Osmar
\end{enumerate}

\vspace{0.4cm}

\begin{longtable}{|P{1.6cm}|P{4.5cm}|P{4.8cm}|P{2.8cm}|}
\hline
\rowcolor{colA}
{\color{white}\textbf{Equipo}} &
{\color{white}\textbf{Actividad}} &
{\color{white}\textbf{Entregas clave}} &
{\color{white}\textbf{Período}} \\
\hline
\endhead

\rowcolor{colC}
Todos &
Análisis de requerimientos, conformación de sub equipos, diseño de la arquitectura
general y definición de las estructuras de datos del sistema. &
Diagrama de clases preliminar y acuerdo sobre las interfaces entre capas. &
Lunes 12/10. \\
\hline

\rowcolor{colD}
1 &
Implementación de los algoritmos de clausura (\texttt{closure}), ir\_a (\texttt{goto})
y construcción de la colección canónica con generación de traza. &
Módulos \texttt{algoritmos\_lr0.py} y \texttt{algoritmo\_lr0.py} funcionales y
verificados con las seis gramáticas de ejemplo. &
Mar.--Vie. 13--16/10. \\
\hline

\rowcolor{colC}
2 &
Implementación de las estructuras de datos (\texttt{Production}, \texttt{Grammar},
\texttt{LR0Item}, \texttt{LR0State}, \texttt{CanonicalCollection}) y del lector de
gramáticas en sus dos formatos. &
Módulos \texttt{modelos\_lr0.py} y \texttt{lector\_gramatica.py} con validaciones
completas y soporte para $\varepsilon$. &
Mar.--Vie. 13--16/10. \\
\hline

\rowcolor{colD}
3 &
Diseño e implementación de la interfaz gráfica: ventana principal, componente de
visualización de estados con scroll y ventana de traza de construcción. &
Módulos de presentación integrados en \texttt{interfaz.py}; interfaz funcional
probada en Windows y macOS. &
Mar.--Vie. 13--16/10. \\
\hline

\end{longtable}

\vspace{0.4cm}

\begin{itemize}
  \item El lunes 12 se dedicó íntegramente al arranque común: análisis de requerimientos,
        conformación de sub equipos y diseño de la arquitectura.
  \item Del martes 13 al viernes 16 cada sub equipo desarrolló en paralelo su módulo
        asignado, con reuniones diarias de sincronización para resolver dependencias.
  \item Los últimos dos días se reservaron para pruebas de integración, corrección de
        errores y redacción del informe de actividades.
\end{itemize}
